\begin{abstract}
Studies of fraud detection on public payment data commonly evaluate models with a random train/test split, and some apply oversampling before the split. In both cases, information from the test set can reach the model during training. We measure how much these practices inflate reported performance by training five standard classifiers (logistic regression, random forest, XGBoost, LightGBM and a small neural network) under random and chronological splits of the IEEE-CIS and ULB credit card datasets, with identical tuning budgets and training-set sizes. On IEEE-CIS, the area under the precision--recall curve is lower by \num{ieee.min.drop.pr} to \num{ieee.max.drop.pr} (\num{ieee.min.rel.pr}--\num{ieee.max.rel.pr}\%) when models are evaluated on later transactions. At a 5\% alert budget, the best model catches \num{ieee.lgbm.chrono.r05} of fraud cases on future transactions, compared with \num{ieee.lgbm.random.r05} under the random split. Controls that score both protocols on the same transactions provide strong evidence that temporal interleaving, rather than a harder test period, explains most of this gap. Under a random split, \num{ieee.key.fraud.random}\% of test fraud cases come from a card with fraud in the training data, compared with \num{ieee.key.fraud.chrono}\% under a chronological split. On the two-day ULB dataset, the same controls attribute most of an apparent drop to a harder final period, and the estimates are too noisy to rank models. Applying SMOTE before splitting is the more serious problem on that dataset, where non-linear models reach a precision--recall AUC of up to \num{ulb.max.leakreal.pr} on genuine test transactions, compared with at most \num{ulb.max.smote.pr} when only training data are oversampled. We provide a checklist for evaluating fraud models and release the code.
\end{abstract}
